% ==================================================
% Telescoping Sums
% Discrete Calculus
% ==================================================

\subsection{Telescoping Sums}
\label{sec:telescoping-sums}

\begin{exposition}
One of the most useful patterns in summation is the phenomenon of
\emph{telescoping}. In a telescoping sum, consecutive terms cancel with
one another, leaving only a small number of terms from the beginning
and end of the expression.
\end{exposition}

\subsection*{A Simple Example}

\begin{example*}
Consider the sum
\[
(2-1) + (3-2) + (4-3) + (5-4).
\]
Expanding gives
\[
2 - 1 + 3 - 2 + 4 - 3 + 5 - 4.
\]
Most terms cancel in adjacent pairs:
\[
\cancel{2} - 1 + \cancel{3} - \cancel{2} + \cancel{4} - \cancel{3} + 5 - \cancel{4}.
\]
After cancellation, only the first and last terms remain: \(5 - 1\).
Thus
\[
(2-1) + (3-2) + (4-3) + (5-4) = 5 - 1.
\]
\end{example*}

\begin{definitionbox}{Definition (Telescoping Sum)}
\begin{definition}[Telescoping Sum]
\label{def:telescoping-sum}
A sum is called a \emph{telescoping sum} if consecutive terms cancel when
the sum is expanded, leaving only a few remaining terms --- typically
the first and last.
\end{definition}
\end{definitionbox}

\begin{remark*}[Standard quantified statement]
A telescoping sum is a sum whose expanded consecutive terms cancel in pairs,
leaving only boundary terms.
\end{remark*}

\begin{remark*}[Interpretation]
The definition names the cancellation pattern used throughout finite
difference and summation arguments.
\end{remark*}

\DefinitionalRoot

\subsection*{General Form}

\begin{exposition}
The general telescoping structure is
\[
\sum_{k=a}^{b-1} (f(k+1) - f(k)).
\]
Writing out the terms explicitly:
\[
\begin{aligned}
& (f(a+1)-f(a)) \\
& + (f(a+2)-f(a+1)) \\
& + (f(a+3)-f(a+2)) \\
& \quad \vdots \\
& + (f(b)-f(b-1)).
\end{aligned}
\]
All intermediate terms cancel in pairs, leaving only \(f(b) - f(a)\).
\end{exposition}

\begin{remark*}[Interpretation]
The name comes from an old collapsible telescope: each section slides
into the next. The intermediate terms in the sum ``slide together''
and vanish, collapsing the entire sum down to its endpoints.
\end{remark*}
